% ===========================================================================
%  2 — O produto
% ===========================================================================
\chapter{O produto}
\label{cap:produto}

\section{O problema}

Crianças e adolescentes vivem dentro de redes sociais, vídeos curtos,
aplicativos de mensagem e assistentes de inteligência artificial --- e a
escola quase nunca consegue falar sobre isso \emph{de dentro}. Falar de
\enquote{fake news} numa aula expositiva é fácil; fazer a turma
\emph{sentir} como um boato viaja, como um feed decide o que cada um vê, ou
por que um vídeo com muitos inscritos não é necessariamente verdadeiro, é
difícil. E o professor da rede pública raramente tem tempo para preparar
esse tipo de aula do zero.

Ao mesmo tempo, as redes públicas precisam cumprir as habilidades da BNCC de
Computação, planejar o ano letivo, cadastrar milhares de pessoas e prestar
contas do que acontece nas escolas --- sem expor dado de criança.

\section{A proposta}

A Plataforma Unooria é uma plataforma educacional vendida para \textbf{redes
públicas de ensino} (secretarias estaduais e municipais). Ela entrega:

\begin{itemize}
  \item \textbf{Missões imersivas para os estudantes}: atividades digitais que
    imitam os aplicativos que eles conhecem (uma rede social de fotos, um
    feed de vídeos curtos, um grupo de mensagens, um assistente de IA, um jogo
    em 3D). O estudante age como agiria no próprio celular; depois, a missão
    revela o que estava acontecendo por trás, e a turma discute.
  \item \textbf{Ferramentas para o professor}: preparar e conduzir a aula,
    gerar plano de aula e planejamento de ensino alinhados à BNCC, e cuidar do
    próprio bem-estar.
  \item \textbf{Gestão da rede}: cadastrar escolas, pessoas e turmas,
    acompanhar a adoção e controlar o que cada rede contratou.
\end{itemize}

\section{Quem compra e quem usa}

O cliente \textbf{não é a escola nem a família}: é a rede. Uma rede contrata a
plataforma e recebe uma instalação própria, separada das instalações de outras
redes. Dentro dela, cem ou mais escolas compartilham a mesma plataforma, e
cada escola só enxerga os próprios dados.

\begin{tabularx}{\linewidth}{@{}L{46mm}X@{}}
\toprule
\headrow \thd{Grandeza, por rede} & \thd{Alvo}\\
\midrule
Estudantes cadastrados & cerca de \num{10000}\\
\zebra Pessoas ao mesmo tempo no horário de pico & \num{2000} a \num{4000}, entre
  7h50 e 8h10 de um dia letivo\\
Escolas & 100 ou mais\\
\zebra Faixa atendida & do 1º ano do Ensino Fundamental à 3ª série do Ensino
  Médio\\
Missões no catálogo alvo & 48 (12 anos $\times$ 4 unidades)\\
\zebra Entrada em produção & novembro de 2026\\
\bottomrule
\end{tabularx}

\section{Os quatro produtos}

O que a rede contrata são \textbf{produtos}. Cada produto aparece como um
cartão no início da área do professor, e a Unooria pode ligar ou desligar cada
um por rede, regional ou escola.

\begin{tabularx}{\linewidth}{@{}L{40mm}L{30mm}X@{}}
\toprule
\headrow \thd{Produto (nome na tela)} & \thd{Chamada} & \thd{O que é}\\
\midrule
\textbf{Entre Redes} & Letramento digital & A trilha de missões: 4 unidades
  por ano, uma missão por unidade, do 1º ano EF à 3ª série EM, mais missões
  \enquote{para qualquer ano}. Hoje: 8 das 48 unidades publicadas, mais a
  Vila do Lume.\\
\zebra \textbf{Gerador de Plano de Aula} & Autoria & Gera um plano anual a
  partir do currículo e do material didático que a rede já adota.\\
\textbf{Planejamento de Ensino} & Currículo e planejamento & Parte da BNCC e
  desce do plano anual até o plano de cada semana, com as datas reais do
  calendário.\\
\zebra \textbf{Como você está?} & Bem-estar & Um acompanhamento curto e
  privado do bem-estar do próprio professor.\\
\bottomrule
\end{tabularx}

\begin{nota}
\textbf{Instanooria, UnooTube e unooGPT não são produtos.} São os nomes dos
aplicativos que as missões imitam (os \enquote{clones}), e vivem todos dentro
do Entre Redes. O estudante nunca vê a palavra \enquote{produto}: ele vê só as
missões da turma dele.
\end{nota}

\section{O caminho de uma aula}
\label{sec:caminho}

A plataforma é um encadeamento de responsabilidades. Cada pessoa faz a sua
parte, e a parte seguinte só acontece se a anterior foi feita. A
Figura~\ref{fig:caminho} mostra o caminho completo de uma missão, da
contratação até o relatório da diretoria.

\begin{figure}[H]
\centering
\fitwidth{%
\begin{tikzpicture}[
  node distance=5mm and 6mm,
  passo/.style={draw=unooBlue, fill=unooBlueWash, rounded corners=2pt,
    minimum width=38mm, minimum height=19mm, align=center, font=\small,
    text width=36mm},
  quem/.style={font=\scriptsize\bfseries, text=unooLavenderSoft},
  liga/.style={-{Stealth[length=2.4mm]}, thick, draw=unooNavy}
]
  \node[passo] (p1) {\textbf{1. Disponibilizar}\\o produto e as missões\\existem nesta rede};
  \node[passo, right=of p1] (p2) {\textbf{2. Cadastrar}\\escolas, pessoas\\e turmas};
  \node[passo, right=of p2] (p3) {\textbf{3. Preparar e liberar}\\a missão para\\a turma};
  \node[passo, below=14mm of p3] (p4) {\textbf{4. Jogar}\\a missão, no ritmo\\da aula};
  \node[passo, left=of p4] (p5) {\textbf{5. Conduzir e ler}\\a turma: debate,\\revelação, resultados};
  \node[passo, left=of p5] (p6) {\textbf{6. Acompanhar}\\a adoção, sem\\identificar criança};

  \node[quem, above=1mm of p1] {Administrador Unooria};
  \node[quem, above=1mm of p2] {Coordenação / Adm. do contrato};
  \node[quem, above=1mm of p3] {Professor};
  \node[quem, below=1mm of p4] {Estudante};
  \node[quem, below=1mm of p5] {Professor};
  \node[quem, below=1mm of p6] {Coordenação regional / Adm. do contrato};

  \draw[liga] (p1) -- (p2);
  \draw[liga] (p2) -- (p3);
  \draw[liga] (p3) -- (p4);
  \draw[liga] (p4) -- (p5);
  \draw[liga] (p5) -- (p6);
\end{tikzpicture}}
\caption[O caminho de uma aula]{O caminho de uma aula. Cada passo depende do
  anterior: sem o produto disponível a missão nem aparece; sem turma o
  estudante não vê missão nenhuma; sem liberação a missão fica com cadeado.}
\label{fig:caminho}
\end{figure}

\begin{atencao}
\textbf{Disponível não é liberado.} São duas perguntas diferentes, feitas por
pessoas diferentes, e confundir as duas é o erro mais caro de suporte:
\begin{itemize}
  \item \textbf{Disponibilidade} --- \enquote{este recurso existe NESTA REDE?}
    Quem responde é só a Unooria, no painel dela. Por padrão, tudo está
    ligado; desligar é exceção (não contratado, ou problema numa versão).
  \item \textbf{Liberação} --- \enquote{esta missão está aberta NESTA TURMA?}
    Quem responde é o professor, na sala. Por padrão, nada está aberto.
\end{itemize}
Um produto desligado tem, para quem está na escola, o mesmo sintoma de um
estudante mal enturmado: a missão não aparece.
\end{atencao}

\section{As pessoas}
\label{sec:pessoas}

São seis papéis com conta, mais o visitante do site. Os nomes abaixo são os
que a interface mostra; nas telas de gestão diz-se sempre
\enquote{Coordenação}, nunca \enquote{gestor}.

\begin{tabularx}{\linewidth}{@{}L{36mm}X X@{}}
\toprule
\headrow \thd{Papel (nome na tela)} & \thd{Quem é} & \thd{O que mais precisa}\\
\midrule
\textbf{Estudante} & Criança ou adolescente da rede, do 1º ano EF à 3ª série
  EM. É titular de dado de menor. & Entrar sem atrito, achar a missão do dia e
  jogar sem perder o que fez.\\
\zebra \textbf{Professor} & Dá aula em uma ou mais turmas, às vezes em mais de
  uma escola. & Preparar a aula rápido, conduzir o momento coletivo, ler a
  turma e planejar o ano.\\
\textbf{Coordenação da escola} & A secretaria ou coordenação pedagógica de uma
  escola. & Cadastrar e manter pessoas e turmas da escola.\\
\zebra \textbf{Coordenação regional} & Diretoria de ensino, polo ou secretaria
  municipal que agrupa escolas. & Ver as escolas da regional, trabalhar numa
  por vez e acompanhar a adoção.\\
\textbf{Administrador do contrato} & A equipe da secretaria que contratou a
  plataforma. & Organizar a rede inteira e ler o agregado, sem entrar na sala
  de aula.\\
\zebra \textbf{Administrador Unooria} & A equipe da Unooria. & Decidir o que
  existe em cada rede, controlar versões e limites de IA, receber
  interessados.\\
\textbf{Visitante} & Alguém de uma rede que ainda não é cliente. & Entender a
  proposta e deixar contato.\\
\bottomrule
\end{tabularx}

\medskip
\textbf{Quem cadastra quem} segue uma cadeia, sempre de cima para baixo:
Administrador Unooria \seta{} Administrador do contrato \seta{} Coordenação
regional \seta{} Coordenação da escola \seta{} Professor e Estudante. Cada
nível cria os que estão abaixo dele, nunca o próprio nível nem um acima.

\section{As nossas promessas}
\label{sec:promessas}

Algumas regras não são detalhe de tela: são o que a plataforma promete à rede,
à escola e à família. Elas valem para todas as missões que existem e para todas
as que vierem, e toda história deste documento respeita todas elas.

\begin{enumerate}
  \item \textbf{A missão só abre na turma em que o professor a liberou.}
    Esconder o cartão não basta: digitar o endereço da missão também não abre.
  \item \textbf{A rede controla o que existe nela --- e desligar nunca apaga.}
    Tudo pode ser desligado por rede, regional ou escola, com motivo e com um
    histórico que a própria rede consegue ler.
  \item \textbf{Nota e correção são da plataforma, nunca do aparelho do
    estudante.} A resposta certa não chega ao tablet antes da hora.
  \item \textbf{Não existe ranking de crianças.} Nenhuma tela de professor ou
    de gestão ordena estudantes por acerto.
  \item \textbf{O relatório não rotula criança.} Ele não afirma o que não mediu:
    \enquote{não teve a oportunidade} nunca é lido como \enquote{escolheu não},
    e um valor com uma única medida não recebe número.
  \item \textbf{Dado de criança é protegido.} Uma escola nunca vê outra; o que
    a criança escreve em texto livre é lido pelo professor dela, não pela
    gestão; os registros de uso não carregam nome.
  \item \textbf{A IA tem limite e tem garantia.} Nenhuma aula depende da IA
    estar no ar; nada gerado chega à turma sem o professor aprovar; a IA nunca
    corrige e nunca recebe dado pessoal; cada pessoa tem um teto diário.
  \item \textbf{O bem-estar do professor é só dele.} Não existe tela de gestão,
    agregado nem exportação desse dado.
  \item \textbf{Tudo em português do Brasil}, com datas no formato brasileiro
    --- exceto dentro dos aplicativos imitados, que falam como o aplicativo
    real (\enquote{Feed}, \enquote{Shorts}, \enquote{Stories}).
\end{enumerate}

\section{Glossário de produto}
\label{sec:glossario}

\begin{tabularx}{\linewidth}{@{}L{34mm}X@{}}
\toprule
\headrow \thd{Termo} & \thd{Significado}\\
\midrule
\textbf{Rede} / contrato & A secretaria cliente. Um contrato é uma rede, com
  instalação própria.\\
\zebra \textbf{Regional} & O nível que agrupa escolas dentro da rede
  (diretoria de ensino, polo, secretaria municipal). Toda escola pertence a
  uma.\\
\textbf{Escola} & A unidade. Tudo que é dado de estudante pertence a uma
  escola.\\
\zebra \textbf{Vínculo} & A ligação de uma pessoa com uma escola, com papel e
  data. Um professor pode ter vários; um estudante tem um ativo por vez.\\
\textbf{Turma} & A sala. O estudante tem uma \textbf{turma de referência}
  (a principal), que é a que as missões usam.\\
\zebra \textbf{Enturmamento} & O ato de pôr pessoas em turmas.\\
\textbf{Ano letivo} / período & Ano, semestre, bimestre ou trimestre. O
  \enquote{8A de 2026} e o \enquote{8A de 2027} são turmas diferentes, o que
  preserva a história.\\
\zebra \textbf{Missão} & O que o estudante joga, visto do ponto de vista da
  aula. É a palavra que professor e estudante leem na tela.\\
\textbf{Unidade} & Cada ano tem 4 unidades; cada unidade tem uma missão.\\
\zebra \textbf{Trilha} & O conjunto das missões do Entre Redes, organizado
  por ano e unidade.\\
\textbf{Clone} / disfarce & A tela que imita um aplicativo real. Dentro dela,
  tudo fala a língua do aplicativo imitado.\\
\zebra \textbf{Revelação} & O momento em que a missão deixa de fingir ser o
  aplicativo e passa a ser a escola falando: a explicação, o \enquote{Meu
  percurso}, o debate.\\
\textbf{Ficha da missão} & A página que o professor lê antes: como a missão
  funciona e o que ela \emph{não} conta ao estudante.\\
\zebra \textbf{Guia da aula} & O roteiro de condução: minuto a minuto,
  perguntas de debate, mediação, avaliação e habilidade da BNCC.\\
\textbf{Liberar} / encerrar & Abrir e fechar a missão para uma turma.\\
\zebra \textbf{Zerar a turma} & Apagar os dados de uma missão numa turma, para
  repeti-la. Não toca outras missões.\\
\textbf{Disponibilidade} & Se um produto ou missão existe numa rede, regional
  ou escola. Quem decide é a Unooria.\\
\zebra \textbf{Limite diário de IA} & O quanto cada pessoa pode usar de IA por
  dia, por perfil (estudante ou professor). Na tela aparece como
  \enquote{Tokens por dia}.\\
\textbf{Perfil portátil} & O documento assinado que o estudante leva ao sair
  da rede, a pedido da família.\\
\zebra \textbf{Primeiro acesso} & A primeira entrada, em que a pessoa troca a
  senha provisória pela dela.\\
\textbf{Equipe} & Quem divide um computador numa missão: dois ou três
  estudantes diante da mesma tela, com o login de um deles, que representa o
  grupo. Vale para uma missão numa turma; não é grupo fixo da turma.\\
\zebra \textbf{Atividade} & Uma tarefa dentro de uma missão, definida pela
  própria missão --- não existe, por ora, uma atividade comum a todas. Hoje só o
  Instanooria tem atividades.\\
\textbf{Rodada} & Cada vez que uma atividade é feita. Reiniciar a atividade
  abre uma rodada nova, e a anterior fica guardada; o resultado é a melhor
  rodada de cada habilidade.\\
\zebra \textbf{Régua pedagógica} / ficha pedagógica & O que uma missão mede:
  quais habilidades da BNCC, em quais momentos, qual pergunta e quanto vale
  cada alternativa. Escrita e validada pela equipe pedagógica.\\
\textbf{Evidência} & O nível que um estudante demonstrou numa habilidade da
  BNCC, numa atividade, numa rodada: atingiu, em processo ou não atingiu. Sem
  registro, é \enquote{Não avaliado}, que nunca conta como não atingiu.\\
\bottomrule
\end{tabularx}
